% Generated draft: requires mathematical and source review.
\documentclass[11pt]{article}
\usepackage[margin=25mm]{geometry}
\usepackage{fontspec}
\setmainfont{Latin Modern Roman}
\usepackage{amsmath,amssymb,mathtools,unicode-math}
\setmathfont{Latin Modern Math}
\usepackage{enumitem,needspace,xurl,hyperref}
\setlength{\parindent}{0pt}
\setlength{\parskip}{5pt}
\setlength{\emergencystretch}{3em}
\newcommand{\N}{\mathbb{N}}
\newcommand{\Z}{\mathbb{Z}}
\newcommand{\Q}{\mathbb{Q}}
\newcommand{\R}{\mathbb{R}}
\newcommand{\C}{\mathbb{C}}
\newcommand{\F}{\mathbb{F}}
\newcommand{\A}{\mathbb{A}}
\newcommand{\PP}{\mathbb{P}}
\newcommand{\E}{\mathbb{E}}
\newcommand{\eps}{\varepsilon}
\newcommand{\dd}{\,\mathrm{d}}
\begin{document}
\section*{UnsolvedMath ID: 800066}
\textbf{Source problem code:} catalogue entry
\subsection{Definitions and mathematical statement}
The supplied record does not specify the class of manifolds \(M\), the contact invariant \(c\), or its pairing. The following definitions explain the usual three-dimensional terminology conditionally; they do not resolve these omissions.

If \(M\) is an oriented smooth three-dimensional manifold, a cooriented contact structure is a plane field \(\xi=\ker\alpha\), where \(\alpha\) is a globally defined one-form and \(\alpha\wedge\dd\alpha\) is nowhere zero. It is positive or negative according to whether \(\alpha\wedge\dd\alpha\) agrees or disagrees with the orientation of \(M\). The source describes the structures as ``(co)oriented'' without specifying the intended orientation data. In this three-dimensional setting, a foliation by surfaces is Reebless if it contains no Reeb component, meaning an embedded solid torus with the standard Reeb foliation, whose boundary is a torus leaf and whose interior leaves are planes.

With the manifold hypotheses and invariant conventions left as intended by the source, let \((\xi_-,\xi_+)\) be a pair of (co)oriented contact structures of opposite signs on \(M\). Assume
\[
\bigl\langle c(\xi_-),c(\xi_+)\bigr\rangle=1.
\]
Must \(M\) admit a Reebless foliation? Equivalently, the requested implication is
\[
\bigl\langle c(\xi_-),c(\xi_+)\bigr\rangle=1
\quad\Longrightarrow\quad
\text{there exists a Reebless foliation on }M.
\]
Here the groups containing the contact invariants, the pairing, and the coefficient system in which \(1\) is interpreted remain unspecified. The source also leaves the dimension, compactness, boundary conditions, and orientation hypotheses on \(M\), and the regularity and transverse orientation requirements for the foliation, unstated. These conventions must be recovered before this becomes a fully specified mathematical statement.
\subsection{Short English statement}
Nariman and Yazdi ask whether a pair of oppositely signed (co)oriented contact structures whose contact invariants pair to \(1\) forces the underlying manifold to admit a Reebless foliation. The supplied record leaves essential manifold and invariant conventions unresolved.
\subsection{Sources}
\begin{itemize}
\item Supplied local UnsolvedMath registry record, permanent ID \(800066\), exact source problem code \texttt{catalogue entry}, titled ``Foliations and Diffeomorphism Groups --- Question 7.'' The registry attributes the problem to Nariman and Yazdi, with proposed year 2025, and labels it open; this status is supplied metadata, not an independently verified assessment.
\item Nariman and Yazdi, \textit{Problem List on Foliations and Diffeomorphism Groups} (2025), Question 7. Supplied source URL: \url{https://arxiv.org/abs/2503.18274}. The provenance records access on 2026-10-16 and extraction method \texttt{arxiv\_latex}. These are supplied provenance claims; the source was not independently consulted for this draft.
\end{itemize}
\end{document}
